% Irodov user-provided PDF, original problems 181--190, printed pp.30--32.
\exproblem{181}{30}
一飞船在无外力下最初以恒速 $v_0$ 飞行。为改变方向，
开启喷气发动机，始终使喷气相对飞船的速度大小为 $u$、
方向垂直于\emph{当时}飞船速度。飞船质量由 $m_0$ 减到 $m$
即停止喷气。求飞行方向的总偏转角。
\exsolution
喷气推力大小 $u\mu$，与飞船瞬时速度垂直，
故 $m\dot v=0$，飞船速率始终为 $v_0$。
推力在速度法线方向产生
\[
m v_0\frac{\mathrm d\alpha}{\mathrm dt}=u\mu,
\qquad \mu=-\frac{\mathrm dm}{\mathrm dt}>0.
\]
于是 $\mathrm d\alpha=-(u/v_0)\,\mathrm dm/m$，
从质量 $m_0$ 积分到 $m$，得到
\[
\boxed{\alpha=\frac{u}{v_0}\ln\frac{m_0}{m}}.
\]
\excheck 角度采用弧度；此结果要求喷气方向始终跟随飞船
当前运动方向转动，不是固定与最初速度垂直。

\exproblem{182}{30}
一辆装沙的小车在光滑水平面上，受与前进方向相同的恒力
$F$。沙从车底孔以恒定质量流率 $\mu>0$ 漏出；
漏出时相对车没有水平速度。最初小车与沙总质量
$m_0$、初速零。求任意 $0\le t<m_0/\mu$ 时的加速度、速度。
\exsolution
漏出的沙在离车一刻具有与车相同的水平速度，
因此没有相对水平喷流推力。车与剩余沙总质量
$m(t)=m_0-\mu t$，水平运动方程
\[
m(t)\frac{\mathrm dv}{\mathrm dt}=F.
\]
直接可得
\[
\boxed{a(t)=\frac{F}{m_0-\mu t}},\qquad
\boxed{v(t)=\int_0^t\frac{F\,\mathrm ds}{m_0-\mu s}
=\frac F\mu\ln\frac{m_0}{m_0-\mu t}}.
\]
\excheck 结果仅在沙尚未漏完的阶段成立；当剩余质量趋零的
数学发散不代表真实小车质量会趋零，应在实际空车质量处终止模型。

\exproblem{183}{30}
一质量为 $m_0$ 的空平板车最初静止，
受沿水平方向的恒力 $F$。静止料斗以恒质量流率
$\mu>0$ 向车中注入沙子（图1.46），沙子落入前
水平速度为零，摩擦忽略。求装料阶段速度与加速度的时间规律。
\exsolution
在固定惯性系中，车与已装入沙的质量为
$m(t)=m_0+\mu t$。
进入系统的沙子带入的水平动量通量为零，
故变质量系统的水平动量平衡为
\[
\frac{\mathrm d}{\mathrm dt}[m(t)v(t)]=F.
\]
由初始总水平动量零，积分得到
$m(t)v(t)=Ft$，于是
\[
\boxed{v(t)=\frac{Ft}{m_0+\mu t}},\qquad
\boxed{a(t)=\frac{\mathrm dv}{\mathrm dt}
=\frac{Fm_0}{(m_0+\mu t)^2}}.
\]
\excheck 此题与上一题不同：正在装入的沙子在地面系
水平静止，必须由车加速，不能错误地套用 $m\dot v=F$。

\exproblem{184}{31}
图1.47中，一根长度为 $\ell$ 的均匀柔软链条置于光滑
水平管内，一段竖直下垂长度为 $h$，其末端恰触桌面。
释放管中另一端 $A$，链条开始从管中滑出，
下垂的链条逐渐在桌面堆积。求 $A$ 刚滑出管口的速率。
\exsolution
设 $x$ 是从管内滑出的累计长度。
因竖直管口至桌面高度固定为 $h$，
已经堆积在桌面上的链长为 $x$，
仍在运动的链长为 $\ell-x$。
设线密度为 $\lambda$，运动链条的重力驱动力
始终为 $\lambda gh$。
在桌面上终止运动的链元携走与当前链条相同的动量，
对运动链段应用含质量流出的动量平衡，通量项抵消
质量随时间变化的相应项，得到
\[
\lambda(\ell-x)\frac{\mathrm dv}{\mathrm dt}
=\lambda gh .
\]
再用 $\dot x=v$，得
\[
v\frac{\mathrm dv}{\mathrm dx}=\frac{gh}{\ell-x}.
\]
从 $v(0)=0$ 积分到 $x=\ell-h$（此时 $A$ 恰到管口）：
\[
\frac{v^2}{2}
=gh\int_0^{\ell-h}\frac{\mathrm dx}{\ell-x}
=gh\ln\frac{\ell}{h}.
\]
所以
\[
\boxed{v=\sqrt{2gh\ln\frac{\ell}{h}}}.
\]
\excheck 不能直接把重力势能损失全部算成末态动能；
链元落桌后非弹性地停住，存在相应动能损失。
需有 $0<h<\ell$，当 $h\to\ell$ 时结果趋零。

\exproblem{185}{31}
某质点相对固定点 $O$ 的角动量矢量随时间变化为
$\bm L(t)=\bm a+\bm b\,t^2$，其中恒定矢量
$\bm a\perp\bm b$。当关于 $O$ 的力矩
$\bm N$ 与角动量夹角为 $45^\circ$ 时，
求力矩大小。
\exsolution
由角动量定理
\[
\bm N=\frac{\mathrm d\bm L}{\mathrm dt}=2t\,\bm b.
\]
设 $a=|\bm a|,\ b=|\bm b|$。
由于 $\bm a\perp\bm b$，
力矩方向沿 $\bm b$（取 $t>0$）；
$\bm L$ 沿 $\bm b$ 的分量大小为 $bt^2$，
垂直分量大小为 $a$。两矢量夹角为 $45^\circ$ 意味着
$bt^2=a$，故
\[
\boxed{|\bm N|=2bt=2\sqrt{ab}}.
\]
\excheck $[\bm b]=[\bm L]/\mathrm s^2$，
所以 $\sqrt{ab}$ 具有力矩量纲；
$t<0$ 时只需用 $|t|$，大小结论不变。

\exproblem{186}{31}
质量为 $m$ 的球以初速度 $v_0$、仰角 $\alpha$
从地面某点抛出，忽略空气阻力。求球对抛出点的
角动量大小随时间变化的规律；并求
$m=0.130\,\mathrm{kg}$、$\alpha=45^\circ$、
$v_0=25\,\mathrm{m/s}$ 时在轨迹顶点的角动量大小。
\exsolution
取抛出点为原点，$x$ 轴水平、$y$ 轴向上：
\[
\bm r=(v_0\cos\alpha\,t,\,
v_0\sin\alpha\,t-\tfrac12gt^2),\quad
\bm v=(v_0\cos\alpha,\,v_0\sin\alpha-gt).
\]
角动量只有垂直于运动平面的分量：
\[
L_z=m(xv_y-yv_x)=-\frac12mgv_0t^2\cos\alpha.
\]
因此
\[
\boxed{|\bm L(t)|=\frac12mgv_0t^2\cos\alpha}.
\]
顶点时间 $t_h=v_0\sin\alpha/g$，
\[
\boxed{L_h=\frac{mv_0^3}{2g}\sin^2\alpha\cos\alpha
\simeq36.6\,\mathrm{kg\,m^2/s}}.
\]
\excheck 取 $g=9.81\,\mathrm{m/s^2}$ 时约 $37$，
与底本答案的两位有效数字相符。角动量矢量方向按
所选坐标为负 $z$ 方向，题目只要求模。

\exproblem{187}{31}
光滑圆片以速率 $v$ 沿光滑水平面滑向固定光滑墙，
碰撞为理想弹性；入射方向与墙面法线夹角为 $\alpha$，
碰撞点为 $O$（图1.48）。
求：（a）相对哪些点碰撞前后角动量守恒；
（b）墙面内、在运动平面上距 $O$ 为 $\ell$ 的点 $O'$
处，圆片角动量变化的模。
\exsolution
碰撞中墙对圆片的冲量仅沿墙面法线。
弹性反射使法线速度分量反号，切线速度分量不变，
故动量增量大小为
\[
|\Delta\bm p|=2mv\cos\alpha,
\]
方向沿墙面法线。关于任意固定点 $Q$ 的角动量变化
是冲量矩
$\Delta\bm L_Q=(\bm r_O-\bm r_Q)\times\Delta\bm p$。
因此当 $Q$ 位于\emph{过 $O$ 的墙面法线直线}上，
力臂与冲量平行，角动量保持不变。
对于题给的 $O'$，其与 $O$ 之间的线段沿墙，
垂直于冲量，于是
\[
\boxed{|\Delta\bm L_{O'}|
=\ell|\Delta\bm p|=2mv\ell\cos\alpha}.
\]
\excheck 虽然关于撞击点 $O$ 的角动量不变，
关于墙上另一点 $O'$ 一般并不不变；
角动量守恒必须说明参考点。

\exproblem{188}{31}
小球质量 $m$，用长为 $\ell$ 的轻绳悬挂于固定点 $O$，
绕竖直轴以恒角速度 $\omega$ 作水平圆周运动（圆锥摆）。
问关于哪个点角动量矢量守恒；
再求关于绳悬点 $O$ 的角动量矢量在半周内的变化量大小。
\exsolution
记小球运动圆半径为 $\rho=\ell\sin\theta$，
圆平面低于 $O$ 的距离为 $h=\ell\cos\theta$。
关于圆心 $C$ 的位置与动量都在水平面内，
角动量始终为沿竖直轴的恒矢量
$\bm L_C=m\omega\rho^2\bm e_z$。
关于 $O$，将位置矢量写为
$\bm r_{O}=\rho\bm e_r-h\bm e_z$，速度为
$\bm v=\omega\rho\bm e_\varphi$，
所以
\[
\bm L_O
=m\omega\rho^2\bm e_z+
m\omega\rho h\,\bm e_r .
\]
一段半周后 $\bm e_r\mapsto-\bm e_r$，
竖直分量不变，故
\[
|\Delta\bm L_O|=2m\omega\rho h.
\]
圆锥摆的动力学关系为
$\omega^2\ell\cos\theta=g$，
最终
\[
\boxed{|\Delta\bm L_O|
=\frac{2mg\ell}{\omega}
\sqrt{1-\left(\frac g{\omega^2\ell}\right)^2}}.
\]
\excheck 只要 $\omega^2\ell\ge g$ 才有非退化圆锥摆；
题目所求“守恒的参考点”为\emph{运动圆的圆心} $C$。

\exproblem{189}{31}
质量 $m$ 的球从高度 $h$、地面系静止释放，
忽略空气阻力。考察以恒定水平速度
$\bm V$ 平移的惯性系，取球的释放点在该系初始为 $O$。
求自由落体期间关于 $O$ 的角动量矢量增量的模。
\exsolution
设地面系竖直向下为 $y$ 轴，
平移系水平方向为 $x$ 轴。
在平移系中球的速度与位置为
\[
\bm v'=(-V,\,gt),\qquad
\bm r'=(-Vt,\,\tfrac12gt^2).
\]
垂直于运动平面的角动量分量：
\[
L_z'=m(x'v_y'-y'v_x')
=m(-Vgt^2+\tfrac12Vgt^2)
=-\tfrac12mVgt^2.
\]
初始为零，落地时间 $t_f=\sqrt{2h/g}$，
从而
\[
\boxed{|\Delta\bm L|=\frac12mVg\,t_f^2=mVh}.
\]
\excheck 使用运动参考系中\emph{固定}于该系的点 $O$；
关于地面系中的释放点，球垂直下落时角动量始终为零，
两者并不矛盾。

\exproblem{190}{31--32}
水平光滑圆盘绕中心固定竖轴以角速度 $\omega$ 匀速转动，
$t=0$ 时，一小圆片从盘中心以相对盘初速
$v_0$ 向外滑行，盘面不施加水平力。
求在转盘参考系中关于盘心的角动量 $L_z(t)$，
并证明其变化由科里奥利力矩引起。
\exsolution
地面惯性系中圆片从原点沿直线以恒速 $v_0$ 运动，
所以 $r(t)=v_0t$。在转盘系中相对速度
$\bm v'=\bm v-\bm\omega\times\bm r$，
其切向分量为 $v_\varphi'=-\omega r$。
关于盘心的相对角动量为
\[
\boxed{L_z'(t)=mr v_\varphi'
=-m\omega v_0^2t^2},
\qquad |L_z'|=m|\omega|v_0^2t^2.
\]
旋转系中离心惯性力沿径向，关于盘心力矩为零；
科里奥利力 $\bm F_{\rm Cor}=-2m\bm\omega\times\bm v'$
沿竖直轴的力矩为
$N_z^{\rm Cor}=-2m\omega r\dot r
=-2m\omega v_0^2t$。
故
\[
\frac{\mathrm dL_z'}{\mathrm dt}
=-2m\omega v_0^2t=N_z^{\rm Cor}.
\]
\excheck 本版书末只写角动量\emph{模}，
而按 $\omega>0$（逆时针）约定的有符号结果为负；
二者一致。
